\documentclass[a4paper]{article}
\usepackage{ctex}
\ctexset{
proofname = \heiti{证明}
}
\usepackage{amsmath, amssymb, amsthm}
\usepackage{moreenum}
\usepackage{mathtools}
\usepackage{url}
\usepackage{bm}
\usepackage{enumitem}
\usepackage{graphicx}
\usepackage{subcaption}
\usepackage{booktabs}
\usepackage[mathcal]{eucal}
\usepackage{array}

% 若本地有课程模板 iidef.sty，则优先使用；
% 若没有，则使用下面的简单备用定义，保证文件尽量可以直接编译。
\IfFileExists{iidef.sty}{
    \usepackage[thehwcnt = 5]{iidef}

    \thecourseinstitute{清华大学}
    \thecoursename{人工智能原理}
    \theterm{2026年春季学期}
    \hwname{作业}
    \slname{\heiti{解}}
}{
    \makeatletter
    \newcommand{\thecourseinstitute}[1]{\def\@courseinstitute{#1}}
    \newcommand{\thecoursename}[1]{\def\@coursename{#1}}
    \newcommand{\theterm}[1]{\def\@term{#1}}
    \newcommand{\hwname}[1]{\def\@hwname{#1}}
    \newcommand{\slname}[1]{\def\@slname{#1}}
    \thecourseinstitute{清华大学}
    \thecoursename{人工智能原理}
    \theterm{2026年春季学期}
    \hwname{作业}
    \slname{\heiti{解}}
    \newcommand{\courseheader}{
        \begin{center}
            {\Large \heiti \@courseinstitute \quad \@coursename}\\[0.5em]
            {\large \@term \quad \@hwname 5 \quad \@slname}
        \end{center}
        \vspace{1em}
    }
    \newcommand{\name}[1]{\noindent #1 \par\vspace{1em}}
    \newenvironment{solution}{\par\noindent\textbf{解：}\par}{\par}
    \makeatother
}

\begin{document}
\courseheader
\name{姓名：杜一郎\qquad 学号：2023011618}

\begin{enumerate}
\setlength{\itemsep}{3\parskip}

\item 两层感知机实现 XNOR
\begin{solution}
\textbf{(a)}

XNOR 的真值表如下：
\[
\begin{array}{ccc}
\toprule
x_1 & x_2 & x_1 \operatorname{XNOR} x_2\\
\midrule
0 & 0 & 1\\
0 & 1 & 0\\
1 & 0 & 0\\
1 & 1 & 1\\
\bottomrule
\end{array}
\]

从真值表可以看出，XNOR 在两种情况下输出为 $1$：一种是 $x_1=x_2=1$，另一种是 $x_1=x_2=0$。因此它可以分解为
\[
x_1 \operatorname{XNOR} x_2
=
(x_1\land x_2)\lor(\lnot x_1\land \lnot x_2).
\]


\textbf{(b)}

令输入向量为
\[
\bm{x}=
\begin{bmatrix}
x_1\\
x_2
\end{bmatrix}.
\]
隐藏层输出为
\[
\bm{h}=f\left(W^{(1)}\bm{x}+\bm{b}^{(1)}\right),
\]
输出层为
\[
\hat{y}=f\left(W^{(2)}\bm{h}+b^{(2)}\right).
\]

取隐藏层权重和偏置为
\[
W^{(1)}
=
\begin{bmatrix}
1 & 1\\
-1 & -1
\end{bmatrix},
\qquad
\bm{b}^{(1)}
=
\begin{bmatrix}
-1.5\\
0.5
\end{bmatrix}.
\]
因此隐藏层两个神经元分别为
\[
h_1=f(x_1+x_2-1.5),
\]
\[
h_2=f(-x_1-x_2+0.5).
\]
其中，$h_1$ 实现 $x_1\land x_2$，$h_2$ 实现 $\lnot x_1\land \lnot x_2$。

输出层取
\[
W^{(2)}
=
\begin{bmatrix}
1 & 1
\end{bmatrix},
\qquad
b^{(2)}=-0.5.
\]
于是
\[
\hat{y}=f(h_1+h_2-0.5).
\]
这相当于对 $h_1$ 和 $h_2$ 做 OR 运算。

\textbf{(c)}

对四种输入逐一计算，得到下表：
\[
\begin{array}{cccccccc}
\toprule
x_1 & x_2 & z_1^{(1)}=x_1+x_2-1.5 & h_1
& z_2^{(1)}=-x_1-x_2+0.5 & h_2
& z^{(2)}=h_1+h_2-0.5 & \hat{y}\\
\midrule
0 & 0 & -1.5 & 0 & 0.5  & 1 & 0.5  & 1\\
0 & 1 & -0.5 & 0 & -0.5 & 0 & -0.5 & 0\\
1 & 0 & -0.5 & 0 & -0.5 & 0 & -0.5 & 0\\
1 & 1 & 0.5  & 1 & -1.5 & 0 & 0.5  & 1\\
\bottomrule
\end{array}
\]

因此，该两层感知机的输出为
\[
\hat{y}=x_1\operatorname{XNOR}x_2,
\]
\end{solution}

\item 两层网络的反向传播梯度
\begin{solution}
\textbf{(a)}

由于 $w_{11}^{(1)}$ 只通过
\[
w_{11}^{(1)}\rightarrow z_1^{(1)}\rightarrow h_1\rightarrow z^{(2)}\rightarrow \hat{y}\rightarrow L
\]
这条路径影响损失函数，所以由链式法则有
\[
\frac{\partial L}{\partial w_{11}^{(1)}}
=
\frac{\partial L}{\partial \hat{y}}
\cdot
\frac{\partial \hat{y}}{\partial z^{(2)}}
\cdot
\frac{\partial z^{(2)}}{\partial h_1}
\cdot
\frac{\partial h_1}{\partial z_1^{(1)}}
\cdot
\frac{\partial z_1^{(1)}}{\partial w_{11}^{(1)}}.
\]

\[
L=\frac{1}{2}(\hat{y}-y)^2,
\]
所以
\[
\frac{\partial L}{\partial \hat{y}}
=
\hat{y}-y.
\]
又因为
\[
\hat{y}=\sigma(z^{(2)}),
\]
并且
\[
\sigma'(z)=\sigma(z)\bigl(1-\sigma(z)\bigr),
\]
所以
\[
\frac{\partial \hat{y}}{\partial z^{(2)}}
=
\hat{y}(1-\hat{y}).
\]
由
\[
z^{(2)}
=
w_1^{(2)}h_1+w_2^{(2)}h_2+b^{(2)}
\]
可得
\[
\frac{\partial z^{(2)}}{\partial h_1}
=
w_1^{(2)}.
\]
由
\[
h_1=\sigma(z_1^{(1)})
\]
可得
\[
\frac{\partial h_1}{\partial z_1^{(1)}}
=
h_1(1-h_1).
\]
最后，由
\[
z_1^{(1)}
=
w_{11}^{(1)}x_1+w_{12}^{(1)}x_2+b_1^{(1)}
\]
可得
\[
\frac{\partial z_1^{(1)}}{\partial w_{11}^{(1)}}
=
x_1.
\]

将这些结果相乘得到
\[
\frac{\partial L}{\partial w_{11}^{(1)}}
=
(\hat{y}-y)\hat{y}(1-\hat{y})w_1^{(2)}h_1(1-h_1)x_1.
\]


\textbf{(b)}
对第一个隐藏层神经元，
\[
z_1^{(1)}
=
0.5\times 1+0.5\times 0+0.5
=
1.
\]
因此
\[
h_1=\sigma(1)=\frac{1}{1+e^{-1}}\approx 0.7311.
\]

对第二个隐藏层神经元，
\[
z_2^{(1)}
=
0.5\times 1+0.5\times 0+0.5
=
1,
\]
所以
\[
h_2=\sigma(1)\approx 0.7311.
\]

再计算输出层：
\[
z^{(2)}
=
0.5h_1+0.5h_2+0.5.
\]
代入 $h_1=h_2\approx 0.7311$，得到
\[
z^{(2)}
\approx
0.5\times 0.7311+0.5\times 0.7311+0.5
=
1.2311.
\]
因此
\[
\hat{y}
=
\sigma(1.2311)
\approx 0.7740.
\]

由上一问的梯度公式，
\[
\frac{\partial L}{\partial w_{11}^{(1)}}
=
(\hat{y}-y)\hat{y}(1-\hat{y})w_1^{(2)}h_1(1-h_1)x_1.
\]
逐项代入：
\[
\hat{y}-y
=
0.7740-1
=
-0.2260,
\]
\[
\hat{y}(1-\hat{y})
=
0.7740(1-0.7740)
\approx 0.1749,
\]
\[
w_1^{(2)}=0.5,
\]
\[
h_1(1-h_1)
=
0.7311(1-0.7311)
\approx 0.1966,
\]
\[
x_1=1.
\]
所以
\[
\frac{\partial L}{\partial w_{11}^{(1)}}
\approx
(-0.2260)\times 0.1749\times 0.5\times 0.1966\times 1.
\]
计算得
\[
\frac{\partial L}{\partial w_{11}^{(1)}}
\approx -0.003886 \approx -0.0039.
\]

\end{solution}

\end{enumerate}

\end{document}
